\section{Baptism: The Gate of Life in Christ}
\label{sec:baptism}

Baptism is the first and indispensable doorway into the Christian life. Through a washing with water in the name of the Father, and of the Son, and of the Holy Spirit, the fallen child of Adam is buried with Christ in His death, cleansed of every stain of original and personal sin, stamped with an indelible priestly seal, and raised as an adopted child of God into the communion of the Church. Every other sacrament presupposes this new birth: without Baptism no one can validly receive Confirmation, the Holy Eucharist, Penance, Anointing of the Sick, Holy Orders, or sacramental Matrimony. For this reason St.~Thomas Aquinas places Baptism first in necessity even while setting the Eucharist first in dignity (\emph{ST}~III, q.~65, a.~3), and the ancient Church named the font at once the Christian's grave and the womb of the Bride.

\subsection{Prefigurations, Institution, and Apostolic Preaching}
\label{subsec:baptism-scripture}

God prepared the sign of baptismal water from the opening sentence of creation, when ``the spirit of God moved over the waters'' (Gen~1:2). When human wickedness filled the earth, the waters of the Deluge swept away the old world of sin while bearing aloft the wooden ark in which eight souls were preserved. St.~Peter reads that event directly as the prophetic pattern of Christian Baptism:
\begin{quotation}
Which had been some time incredulous, when they waited for the patience of God in the days of Noe, when the ark was a building: wherein a few, that is, eight souls, were saved by water. Whereunto baptism, being of the like form, now saveth you also: not the putting away of the filth of the flesh, but, the examination of a good conscience towards God by the resurrection of Jesus Christ. (1~Pet~3:20--21)
\end{quotation}
The same double movement of judgment and deliverance unfolded at the Red Sea, where Israel walked dry-shod between walls of water while Pharaoh's pursuing host was drowned (Exod~14:21--22). St.~Paul explicitly names that passage a figure of the font: ``our fathers were all under the cloud: and all passed through the sea. And all in Moses were baptized, in the cloud and in the sea'' (1~Cor~10:1--2). Later, Israel crossed the Jordan into the promised land, and the prophet Ezechiel announced the inward purification no Levitical washing could achieve: ``And I will pour upon you clean water, and you shall be cleansed from all your filthiness, and I will cleanse you from all your idols. And I will give you a new heart, and put a new spirit within you'' (Ezech~36:25--26).

These preparatory figures reached fulfillment when the Eternal Word stepped down into the Jordan to be baptized by John. Sinless Himself, Jesus entered the stream not to wash away any guilt of His own, but ``to fulfil all justice'' (Matt~3:15), standing in solidarity with sinners and sanctifying the waters by contact with His sacred flesh. As He rose from the river, the Trinity manifested itself: the heavens opened, the Spirit of God descended like a dove upon Him, and the Father's voice proclaimed, ``This is my beloved Son, in whom I am well pleased'' (Matt~3:16--17). John's baptism was a preparatory washing of repentance; Christ's descent into the Jordan foreshadowed the baptism of His Cross and inaugurated the sanctification of water for the New Law (\emph{ST}~III, q.~66, a.~2).

In His nocturnal discourse with Nicodemus, the Lord stated the necessity and supernatural nature of this new birth plainly: ``Amen, amen, I say to thee, unless a man be born again of water and the Holy Ghost, he cannot enter into the kingdom of God. That which is born of the flesh is flesh: and that which is born of the Spirit is spirit'' (John~3:5--6). On Calvary, the lance opened the side of the dead Redeemer, and ``immediately there came out blood and water'' (John~19:34). St.~John returns to that pierced side in his first epistle, confessing that ``This is he that came by water and blood, Jesus Christ: not by water only but by water and blood,'' and joining ``the spirit and the water and the blood'' on earth (1~John~5:6, 8; in the Clementine Vulgate and Douay-Rheims, the heavenly triad stands at v.~7 before the earthly triad of v.~8). After conquering death, the risen Christ promulgated the universal mandate of the sacrament:
\begin{quotation}
All power is given to me in heaven and in earth. Going therefore, teach ye all nations: baptizing them in the name of the Father and of the Son and of the Holy Ghost. Teaching them to observe all things whatsoever I have commanded you. And behold I am with you all days, even to the consummation of the world. (Matt~28:18--20)
\end{quotation}
St.~Mark preserves the solemn sanction attached to that charge: ``He that believeth and is baptized shall be saved: but he that believeth not shall be condemned'' (Mark~16:16).

From Pentecost morning onward, the Apostles carried out this mission without delay. When the crowd in Jerusalem asked what to do, St.~Peter answered: ``Do penance: and be baptized every one of you in the name of Jesus Christ, for the remission of your sins. And you shall receive the gift of the Holy Ghost'' (Acts~2:38). The expression ``in the name of Jesus Christ'' denotes baptism conferred by Christ's authority and unto His saving faith, never an omission of the Triune formula commanded in Matt~28:19 (\emph{ST}~III, q.~66, a.~6).

St.~Paul anchors every baptismal effect in union with Christ's death and Resurrection:
\begin{quotation}
Know you not that all we who are baptized in Christ Jesus are baptized in his death? For we are buried together with him by baptism into death: that, as Christ is risen from the dead by the glory of the Father, so we also may walk in newness of life. For if we have been planted together in the likeness of his death, we shall be also in the likeness of his resurrection. (Rom~6:3--5)
\end{quotation}
Across the Pauline epistles, converging images illuminate this transformation. Baptism clothes the believer in Christ Himself: ``For as many of you as have been baptized in Christ have put on Christ,'' overcoming every division of Jew and Greek, slave and free (Gal~3:27--28). It incorporates the neophyte into the Mystical Body: ``For in one Spirit were we all baptized into one body'' (1~Cor~12:13). Christ cleanses His Bride ``by the laver of water in the word of life'' (\latin{mundans lavacro aquae in verbo vitae}, Eph~5:26). It stands as an unrepeatable bond of Catholic unity: ``One Lord, one faith, one baptism'' (Eph~4:5). And God saves us not by our own works of justice, but ``by the laver of regeneration and renovation of the Holy Ghost'' (\latin{per lavacrum regenerationis et renovationis Spiritus Sancti}, Titus~3:5).

\subsection{The Patristic Teaching: Illumination, Grave, and Mother}
\label{subsec:baptism-fathers}

The Fathers of the Church beheld Baptism as a divine intervention in which the Holy Trinity descends to recreate a human soul. Writing around the middle of the second century in Rome, St.~Justin Martyr offers the earliest detailed account of the rite outside the New Testament. After describing how converts convinced of the truth promise to live accordingly and join the faithful in prayer and fasting, he describes the washing itself:
\begin{quotation}
Then they are brought by us where there is water, and are regenerated in the same manner in which we were ourselves regenerated. For, in the name of God, the Father and Lord of the universe, and of our Saviour Jesus Christ, and of the Holy Spirit, they then receive the washing with water. For Christ also said, ``Except ye be born again, ye shall not enter into the kingdom of heaven.'' \dots\ And this washing is called illumination, because they who learn these things are illuminated in their understandings. (St.~Justin Martyr, \emph{First Apology}~61)
\end{quotation}
Already in second-century Rome, regeneration (\latin{anagennesis}), illumination (\latin{photismos}), catechesis, repentance, natural water, and the explicit Trinitarian invocation stand seamlessly united.

Two centuries later in Jerusalem, St.~Cyril taught Easter candidates that the preparatory scriptural exorcisms refine the soul like a goldsmith's fire (\emph{Procatechesis}~9), and hailed the approaching bath in a cascade of titles:
\begin{quotation}
Great is the Baptism that lies before you: a ransom to captives; a remission of offences; a death of sin; a new-birth of the soul; a garment of light; a holy indissoluble seal; a chariot to heaven; the delight of Paradise; a welcome into the kingdom; the gift of adoption! (St.~Cyril of Jerusalem, \emph{Procatechesis}~16)
\end{quotation}
In his third \emph{Catechetical Lecture}, St.~Cyril explains how bodily washing reaches an immaterial spirit: ``the simple water having received the invocation of the Holy Ghost, and of Christ, and of the Father, acquires a new power of holiness'' (\emph{Catechetical Lecture}~III.3). Because man is soul and body, the purification matches both parts of his nature: ``the water cleanses the body, and the Spirit seals the soul'' (\emph{Catechetical Lecture}~III.4).

In the \emph{Mystagogical Catecheses} preached during Easter Week to the newly baptized (attributed traditionally to St.~Cyril of Jerusalem, though modern scholarship also weighs authorship by his successor John~II of Jerusalem), the preacher recalls the threefold descent into the pool at the Holy Sepulchre in an astonishing image:
\begin{quotation}
And each of you was asked, whether he believed in the name of the Father, and of the Son, and of the Holy Ghost, and ye made that saving confession, and descended three times into the water, and ascended again; here also hinting by a symbol at the three days burial of Christ. \dots\ And at the self-same moment ye were both dying and being born; and that Water of salvation was at once your grave and your mother. (\emph{Mystagogical Catechesis}~II.4)
\end{quotation}
Though the neophytes did not undergo physical crucifixion, through sacramental participation they entered the real fruit of Calvary: ``our imitation was in a figure, and our salvation in reality'' (\emph{Mystagogical Catechesis}~II.5).

In Milan, St.~Ambrose taught his Easter neophytes in \emph{De mysteriis} to see past the humble baptistery with eyes of faith: ``Water, certainly, but not water alone; you saw the deacons ministering there, and the bishop asking questions and hallowing. \dots\ Believe, then, that the presence of the Godhead is there'' (\emph{De mysteriis}~III.8). Water without the preaching of the Cross avails nothing (\emph{De mysteriis}~III.14); pressing the unity of water, blood, and Spirit from 1~John~5:8, St.~Ambrose insists that every component of Christ's institution is indispensable:
\begin{quotation}
Therefore read that the three witnesses in baptism, the water, the blood, and the Spirit, are one, for if you take away one of these, the Sacrament of Baptism does not exist. For what is water without the cross of Christ? A common element, without any sacramental effect. Nor, again, is there the Sacrament of Regeneration without water \dots\ unless he be baptized in the Name of the Father, and of the Son, and of the Holy Spirit, he cannot receive remission of sins nor gain the gift of spiritual grace. (St.~Ambrose, \emph{De mysteriis}~IV.20)
\end{quotation}

St.~Augustine brought this heritage to speculative precision against the Donatists and Pelagians. Against Donatist re-baptism, he insisted that Christ Himself is always the true Baptizer regardless of the minister's moral state. In Tractate~80.3 on the Gospel of John, St.~Augustine defines the union of element and spoken form and applies it directly to infant baptism:
\begin{quotation}
Take away the word, and the water is neither more nor less than water. The word is added to the element, and there results the Sacrament, as if itself also a kind of visible word. \dots\ This word of faith possesses such virtue in the Church of God, that through the medium of him who in faith presents, and blesses, and sprinkles it, He cleanseth even the tiny infant, although itself unable as yet with the heart to believe unto righteousness, and to make confession with the mouth unto salvation. (St.~Augustine, \emph{In Ioann. Evang.}~Tr.~80.3)
\end{quotation}
As noted in the opening chapter, St.~Augustine dwells in Tractate~120.2 on St.~John's exact verb in John~19:34---``opened'' rather than merely wounded---contemplating Christ's pierced side as the gate of life from which Baptism and the Eucharist issue to form the Church, His Bride, as Eve was formed from the sleeping Adam.

At the close of the patristic age in the East, St.~John of Damascus sums up the Greek Fathers in \emph{De fide orthodoxa}~IV.9: because our Lord died once for all, ``we also must be baptized once for all, and baptized according to the Word of the Lord, In the Name of the Father, and of the Son, and of the Holy Spirit.'' From His immaculate side welled forth ``the fountain of remission \dots\ water for our regeneration, and the washing away of sin and corruption; and blood to drink as the hostage of life eternal.'' While the remission of sins is granted alike to all, ``the grace of the Spirit is proportional to the faith and previous purification'' of those possessing reason.

\subsection{Sacramental Constitution: Matter, Form, Minister, and Subject}
\label{subsec:baptism-constitution}

Adopting Peter Lombard's concise definition---``Baptism is the outward washing of the body done together with the prescribed form of words'' (\emph{ST}~III, q.~66, a.~1)---St.~Thomas Aquinas, the Fourth Lateran Council (\emph{Firmiter credimus}, 1215, DS~802), the Council of Florence (\emph{Exsultate Deo}, 1439), the Council of Trent (Session~VII), the \emph{Roman Catechism} (Part~II), and the modern Codes of Canon Law specify every requirement for valid and licit celebration.

\paragraph{Matter.}
The remote matter of Baptism is true, natural water (\emph{ST}~III, q.~66, a.~3; CIC~1983, c.~849; CCEO~1990, c.~675~\S1). Liquids not recognized as water---such as wine, milk, oil, juice, or broth---are invalid even in emergencies. St.~Thomas gives four fitting reasons for Christ's choice of water (\emph{ST}~III, q.~66, a.~3): first, water's natural kinship with living generation answers to spiritual rebirth; second, its physical properties mirror baptismal effects, since moistness cleanses sin, coolness moderates the heat of concupiscence, and transparency receives light as the sacrament of faith; third, dipping into water and rising again signify Christ's burial and Resurrection; and fourth, water is available everywhere so no person is barred from salvation by scarcity. Outside danger of death, liturgical law prescribes water blessed at the Easter Vigil or during the rite itself, though any natural water suffices for validity.

The proximate matter is the real outward washing (\latin{ablutio}) of the body with that water. Because washing requires water to touch and flow over the skin (ordinarily over the head), Baptism is validly conferred by immersion, by pouring (\latin{infusio}), or by flowing sprinkling (\latin{aspersio}; \emph{ST}~III, q.~66, a.~7; CIC~1983, c.~854). Triple immersion vividly expresses burial and rising (Rom~6:4) and remains standard across Eastern Catholic Churches, while infusion is common lawful practice in the Latin Church.

\paragraph{Form.}
The essential form must simultaneously express the washing act, the person baptized, and the explicit invocation of the Three Persons of the Blessed Trinity commanded in Matt~28:19 (\emph{ST}~III, q.~66, a.~5). Catholic communion receives two apostolic forms of this single Trinitarian invocation (CCC~1240):
\begin{itemize}
  \item \textbf{Latin active form:} ``N., I baptize you in the name of the Father, and of the Son, and of the Holy Spirit'' (\latin{Ego te baptizo in nomine Patris, et Filii, et Spiritus Sancti}).
  \item \textbf{Eastern passive form:} ``The servant of God, N., is baptized in the name of the Father, and of the Son, and of the Holy Spirit'' (with minor traditional phrasing across the five Eastern liturgical families).
\end{itemize}
Both express the same truth: the Latin singular ``I'' manifests Christ Himself acting personally through His minister, while the Eastern passive voice underscores that the Trinity alone sanctifies the candidate. The washing and words must be joined morally in one action by the same minister (\emph{Roman Catechism}, Part~II). Substituting non-Trinitarian titles (such as ``Creator, Redeemer, and Sanctifier'') invalidates the sacrament. Likewise, in its 2020 \emph{Responsum ad dubium} (6~August 2020, approved by Pope Francis), the Congregation for the Doctrine of the Faith confirmed that altering the Latin singular ``I baptize you'' to the communal plural ``We baptize you'' invalidates the rite: the assembled community does not baptize in its own name, and deliberately changing the received form breaks the sign by which Christ acts through the minister---a principle reaffirmed across all seven sacraments in the Dicastery for the Doctrine of the Faith's Note \emph{Gestis Verbisque} (2~February 2024).

\paragraph{Minister and Intention.}
Christ is always the principal minister. In universal canon law, the ordinary ministers of Baptism in the Latin Church are the bishop, priest, and deacon (CIC~1983, c.~861~\S1); in the Eastern Catholic Churches, the ordinary minister is the priest (normally the pastor), with the deacon acting only in necessity or per particular law (CCEO~1990, c.~677~\S\S1--2), reflecting the Eastern union of Baptism and presbyteral Chrismation.

Because Baptism is necessary for salvation, divine mercy extends extraordinary ministry to every human being in danger of death. The Fourth Lateran Council defined in \emph{Firmiter credimus} (1215, DS~802) that Baptism consecrated in water with the invocation of the undivided Trinity avails unto salvation for infants and adults when rightly conferred in the Church's form by anyone whatsoever (\latin{a quocumque rite collatum}). St.~Thomas explains that in necessity even a layperson, a heretic, or an unbaptized person confers Baptism validly, provided he applies natural water, pronounces the Trinitarian form, and holds the minimal right intention ``to do what the Church does'' (\emph{ST}~III, q.~67, aa.~3, 5; Trent, Sess.~VII, can.~4 on Baptism, can.~11 on Sacraments in General; CIC~1983, c.~861~\S2; CCEO~1990, c.~677~\S2).

\paragraph{Capable Subject.}
Every living unbaptized human person---and only such a person---is capable of receiving Baptism (\emph{ST}~III, q.~68, a.~1; CIC~1983, c.~864; CCEO~1990, c.~679). One already validly baptized can never be re-baptized; where serious doubt remains after inquiry, Baptism is administered conditionally (CIC~1983, c.~869; CCEO~1990, c.~672~\S2). An adult with the use of reason needs free intention for validity, and faith with genuine repentance for fruitful grace (\emph{ST}~III, q.~68, aa.~4, 7; CIC~1983, c.~865~\S1). Infants have been baptized from apostolic times because they inherit original sin from Adam and need rebirth in Christ (\emph{ST}~III, q.~68, a.~9; Lateran~IV, \emph{Firmiter}; Trent, Sess.~V, can.~4; Sess.~VII, cans.~12--13 on Baptism), borne to the font in the faith of Mother Church through parents and sponsors (CIC~1983, cc.~867--868, 872--874; CCEO~1990, cc.~681, 684--686). Preparatory exorcisms, renunciation of Satan, profession of the Creed, and pre-baptismal anointing signify liberation from darkness without forming part of the bare minimum required for emergency validity (\emph{ST}~III, q.~71, aa.~2--4).

\begin{fourstable}{Constitutive Dimension}{Metaphysical and Ritual Core}{Latin Church (CIC~1983)}{Eastern Churches (CCEO~1990)}
Matter & Natural water applied by real washing (immersion, infusion, or flowing aspersion) & Blessed water ordinarily; infusion or immersion (cc.~849, 853--854) & Blessed water ordinarily; triple immersion traditional (c.~675~\S1; liturgical books)\\
Form & Explicit Trinitarian invocation joined to washing (Matt~28:19; CCC~1240) & Active: ``N., I baptize you in the name of the Father, and of the Son, and of the Holy Spirit'' & Passive: ``The servant of God, N., is baptized in the name of the Father, and of the Son, and of the Holy Spirit''\\
Minister & Christ principally; human minister with intention to do what the Church does & Ordinary: bishop, priest, deacon (c.~861~\S1); emergency: anyone (c.~861~\S2) & Ordinary: priest, normally the pastor (c.~677~\S1); emergency: deacon or any person (c.~677~\S2)\\
Subject and Sequence & Any living unbaptized human person; adults with faith and repentance & Infants baptized promptly (c.~867); adults receive B--C--E together (c.~866) & Infants and adults ordinarily receive unified B--C--E in one liturgy (cc.~695~\S1, 697, 710)\\
\end{fourstable}

\subsection{Sacramental Effects, Indelible Character, Change, and Order of Ends}
\label{subsec:baptism-effects}

Catholic theology distinguishes three nested layers in Baptism (\emph{ST}~III, q.~66, a.~1):
\begin{itemize}
  \item \textbf{\latin{Sacramentum tantum}} (the outward sign alone): the visible washing with water joined to the Trinitarian words.
  \item \textbf{\latin{Res et sacramentum}} (the intermediate reality and sign): the indelible \emph{baptismal character} impressed upon the soul.
  \item \textbf{\latin{Res tantum}} (the ultimate proper grace): inward justification---complete remission of sin and punishment, sanctifying grace, adoptive filiation, the theological and moral virtues, the gifts of the Holy Spirit, and incorporation into the Mystical Body of Christ.
\end{itemize}

\paragraph{The Baptismal Character.}
Whenever Baptism is validly conferred, even if an insincere adult places an obstacle (\latin{obex}) in the way of grace, it immediately imprints an indelible character upon the intellective faculty of the soul (\emph{ST}~III, q.~63, aa.~1--6; Trent, Sess.~VII, can.~9 on Sacraments in General). This spiritual potency configures the believer to Christ the Priest, enabling valid reception of the other sacraments and participation in the common priesthood of the faithful (\emph{Lumen Gentium}~10--11). Because Christ's priesthood is eternal, no sin can erase the baptismal character; Baptism can never be repeated (\emph{ST}~III, q.~66, a.~9; Eph~4:5).

\paragraph{Remission of Sin and New Life in Grace.}
In every unobstructed recipient, Baptism completely blots out original sin and every prior actual mortal and venial sin (\emph{ST}~III, q.~69, a.~1; Trent, Sess.~V, can.~5), together with the entire debt of eternal and temporal punishment (\latin{poena}) due to past sins (\emph{ST}~III, q.~69, a.~2). Bodily mortality and the disordered inclination called concupiscence (\latin{fomes peccati}) remain, yet the Council of Trent (Session~V, can.~5) defines that concupiscence in the baptized is not truly sin, but arises from sin and inclines toward sin as an arena for spiritual combat (\emph{ST}~III, q.~69, a.~3).

Positively, the \latin{res tantum} of Baptism infuses habitual sanctifying grace---a created participation in the divine nature---together with faith, hope, charity, the infused moral virtues, the sevenfold gifts of the Spirit, and the indwelling Trinity (\emph{ST}~III, q.~69, aa.~4--5). The baptized becomes an adopted child of God, an heir of heaven (Titus~3:7), and a living member of the Church (1~Cor~12:13). Where an adult receives Baptism validly yet insincerely (\latin{cum fictione}), the character is imprinted while grace is blocked; once sincere penance removes that obstacle, baptismal grace takes effect through \emph{reviviscence} (\emph{ST}~III, q.~69, a.~10).

\paragraph{The Metaphysical Change Effected.}
Unlike the Holy Eucharist, where the entire substance of bread and wine changes into Christ's Body and Blood, Baptism effects \textbf{no substantial change} in the recipient. The baptized person remains identically the same human creature in body and rational soul. Rather, Baptism effects a real \textbf{accidental formal elevation} of the soul---through both the indelible character (an accidental spiritual potency in the intellect) and sanctifying grace (an accidental supernatural quality elevating the soul's essence to participate in divine life)---joined to a permanent \textbf{ecclesial relation} incorporating the believer into the Mystical Body of Christ. Grace perfects nature from within while keeping creaturehood wholly intact.

\paragraph{The Order of Ends.}
The ends of Baptism stand in clear teleological order:
\begin{enumerate}
  \item \textbf{Primary proper end:} spiritual regeneration into the life of Christ and complete purification from original and personal sin.
  \item \textbf{Intrinsic secondary ends:} visible incorporation into the Church, conferral of the priestly character enabling Christian worship and valid reception of the remaining six sacraments, and establishment of a real foundational bond of unity among all validly baptized Christians (\emph{Unitatis Redintegratio}~22).
  \item \textbf{Contingent ends:} the spiritual consolation of parents, sponsors, and the local parish celebrating new life in Christ, and the public testimony of an adult convert's turning to the Gospel.
  \item \textbf{Ultimate end:} the glory of the Triune God and eternal beatitude in heaven (John~3:5; Titus~3:7).
\end{enumerate}

\subsection{Necessity, Desire, Martyrdom, and the Bond of Unity}
\label{subsec:baptism-necessity}

Christ's words to Nicodemus (John~3:5) and His parting command (Mark~16:16) establish that Baptism is necessary for salvation by both divine precept and intrinsic means (\emph{necessitate medii}; Trent, Sess.~VII, can.~5 on Baptism). Yet God, who charged His Church to administer the sacraments, did not restrict His own sovereign mercy to them. As St.~Thomas teaches, ``God, Whose power is not tied to visible sacraments,'' can bestow the inward grace of Baptism (\latin{res sacramenti}) without outward water when a person is prevented from reaching the font without fault (\emph{ST}~III, q.~68, a.~2). Following the Fathers, St.~Thomas distinguishes three modes of Baptism (\emph{ST}~III, q.~66, a.~11):

\begin{threestable}{Mode of Baptism}{Nature and Basis}{Sacramental Reality Conferred}
Baptism of Water (\latin{baptismus fluminis}) & Ordinary sacramental washing with natural water and Trinitarian form instituted by Christ (John~3:5; Matt~28:19) & Imprints the indelible baptismal character; remits all sin and punishment; infuses grace and incorporates visibly into the Church\\
Baptism of Blood (\latin{baptismus sanguinis}) & Suffering martyrdom for Christ or Christian virtue by an unbaptized person before reaching the font (\emph{ST}~III, q.~66, a.~11) & Imprints no sacramental character, yet confers full justification and complete remission of sin and punishment by union with Christ's Passion\\
Baptism of Desire (\latin{baptismus flaminis}) & Explicit desire (\latin{votum}) of a catechumen prevented by death, or implicit desire in supernatural faith and perfect charity (\emph{ST}~III, q.~68, a.~2; Trent, Sess.~VI) & Imprints no sacramental character, yet imparts sanctifying grace and remission of guilt; never dispenses from water when opportunity remains\\
\end{threestable}

Neither martyrdom nor desire weakens the Church's missionary mandate to baptize all nations, since only sacramental Baptism imprints the character, incorporates a believer visibly into liturgical communion, and opens access to Confirmation and the Eucharist. Regarding infants who die without Baptism, the Church entrusts them in liturgical prayer and hope to the boundless mercy of God, who wills all men to be saved (1~Tim~2:4; CCC~1261), while upholding the pastoral duty to baptize infants promptly (CIC~1983, c.~867).

Finally, because valid Baptism imprints Christ's character wherever water and the Trinitarian form are used with the Church's intention, the Second Vatican Council teaches in \emph{Unitatis Redintegratio}~22 that Baptism creates a genuine sacramental bond of unity linking all reborn through it. Yet Baptism by itself is a beginning, intrinsically ordered toward full profession of apostolic faith, full ecclesial incorporation, and full participation in the Eucharistic table. Born at the font, the Christian moves directly toward the strengthening chrism of Confirmation and the living Bread of the altar.
